\clearpage
\thispagestyle{plain}
\phantomsection
\addcontentsline{toc}{chapter}{ABSTRAK}

{\centering\fontsize{12}{14}\selectfont\bfseries ABSTRAK\par}

\vspace{0.2cm}

\noindent
\textbf{\MakeUppercase{\JudulPA}}

\vspace{0.2cm}

\noindent
\textbf{Nama\hspace{1.45cm}: \NamaMahasiswa}\\
\textbf{NRP\hspace{1.65cm}: \NRP}\\
\textbf{Pembimbing\hspace{0.2cm}: \PembimbingSatu}\\
\textbf{\phantom{Pembimbing\hspace{0.2cm}: }\PembimbingDua}

\vspace{0.2cm}

Penyakit kardiovaskular merupakan penyebab utama kematian di dunia, dan aritmia sering terjadi tanpa gejala sehingga monitoring elektrokardiogram (ECG) secara kontinu diperlukan untuk deteksi dini. Sistem monitoring yang ada umumnya berarsitektur \textit{cloud-centric} yang berlatensi tinggi dan bergantung pada koneksi internet. Model deteksi aritmia juga jarang diuji pada batasan mikrokontroler, dan \textit{pipeline} data dari \textit{edge} ke \textit{cloud} belum tahan terhadap gangguan jaringan. Penelitian ini bertujuan membangun sistem monitoring detak jantung single-lead berbasis \textit{edge computing} dengan pendekatan TinyML. Preprocessing sinyal ECG dan klasifikasi aritmia dijalankan secara \textit{on-device} pada ESP32-S3, dan hasilnya dikirim melalui \textit{pipeline} MQTT yang dirancang tahan terhadap putusnya koneksi. Filter Butterworth dan deteksi R-\textit{peak} Pan-Tompkins dijalankan di mikrokontroler, dan model CNN \textit{hybrid} dengan cabang morfologi dan cabang ritme dikuantisasi menjadi INT8. Detak yang belum terkirim ditampung pada bufer lokal di PSRAM. Model \textit{float32} pada DS2 MIT-BIH dengan 22 pasien yang tidak terlihat saat pelatihan mencapai \textit{recall} 0,6996, F1 0,6594, dan AUC 0,9334. Kuantisasi INT8 menghasilkan model 22,94~KB dengan selisih \textit{recall} $+$0,0108, dan inferensinya di ESP32-S3 memakan 26,6~ms per detak. Preprocessing pada kode C terbukti setara dengan Python dengan beban 3,0\% anggaran waktu tanpa sampel yang hilang. Keempat rekaman sinyal tubuh resmi dari satu sesi perekaman memenuhi kriteria kualitas, dan gerbang kalibrasi berbasis kontak elektroda membuka pemantauan sekitar 60 detik setelah kontak baik. Dengan sinyal \textit{replay} dari rekaman MIT-BIH, \textit{pipeline} MQTT mencapai \textit{packet delivery success rate} 99,2\% sampai 100,0\% pada tiga ulangan dengan median latensi \textit{end-to-end} 939 sampai 1.108~ms. Pembekuan \textit{broker} selama 24 detik tidak menghilangkan satu detak pun. Tujuan preprocessing dan \textit{pipeline} tercapai, sedangkan tujuan model INT8 tercapai kecuali target konsumsi daya yang belum diukur.
% sumber: bab4 tab:hasil-ds2 (DS2 22 pasien, recall 0,6996, F1 0,6594, AUC 0,9334), tab:delta-int8 (22,94 KB, selisih recall +0,0108), tab:sumber-daya (26,6 ms per detak), tab:beban-komputasi (3,0%, sampel hilang 0), tab:kualitas-rekaman (4 dari 4 rekaman resmi Berhasil, satu sesi), butir Gerbang Kalibrasi (sekitar 60 detik), tab:target-pipeline (PDSR 100,0% dan 99,2%, median 939 dan 1.108 ms, dua ulangan, sumber sinyal replay), butir Putus Koneksi (beku 24 detik, hilang 0), tab:ketercapaian (preprocessing dan pipeline Tercapai, model Sebagian tercapai karena daya)

\vspace{0.25cm}
\noindent
\textbf{Kata kunci:} ECG, IoT, Edge Computing, TinyML, ESP32-S3, Aritmia, MQTT

\clearpage
